\section{Methodology}
\label{sec:method}

\subsection{Dataset Construction}
\label{sec:data}

We construct four purpose-built experimental datasets to enable precise behavioral and
mechanistic analyses.  All datasets use five tool categories: weather forecasting,
mathematical calculation, web search, code execution, and translation.  Every example
presents a query alongside four candidate tools, each described with a realistic name and
description drawn from real API documentation (OpenWeatherMap, WolframAlpha, Google
APIs, etc.).  Queries are designed to be semantically clear but not to contain exact tool
names, preventing trivial lexical matching.

\para{Main tool selection dataset (55 examples).}
The primary dataset contains 40 same-category examples and 5 mixed-category examples
(10 queries $\times$ 5 categories, plus 5 cross-category scenarios).  Same-category
examples present four tools from the same functional domain, requiring fine-grained
semantic matching; mixed-category examples present one tool per domain, testing coarse
routing.  The position of the ground-truth tool is randomized uniformly across the four
slots.

\para{Positional bias dataset (40 variants).}
We create four cyclic rotations of the tool list for each of 10 base examples, yielding
40 variants total.  The same query is paired with the same four tools in four different
orderings, allowing us to isolate order effects from content effects.

\para{Description perturbation dataset (7 examples).}
We test three conditions for each query: \emph{original} (real descriptions),
\emph{generic} (all descriptions replaced with ``An API service that processes requests
and returns data''), and \emph{adversarial} (descriptions swapped between tools, so the
semantically correct tool carries a misleading description).  Three distinct query types
are covered (weather, calculation, translation), yielding 7 prompts.

\para{Probing dataset (100 examples).}
We construct 100 tool-selection prompts with 25 examples per category (4 categories used
in probing) to obtain balanced hidden-state representations for the linear probing
experiments.

\subsection{Behavioral Analysis with \gptmini}
\label{sec:llm_expts}

We collect tool selections from \gptmini via the OpenRouter API with temperature~$= 0$
(deterministic) and a 50-token output budget.  Each example is run three times to measure
consistency; we use the majority-vote selection as ground truth.  All API calls follow a
structured prompt format that presents the query and asks the model to name the most
appropriate tool.

\subsection{Semantic Similarity Analysis}
\label{sec:semsim}

We embed each query and each tool description using \minilm (22M parameters,
sentence-transformers library).  For each example we rank the four tools by cosine
similarity to the query and record whether the highest-ranked tool matches the \gptmini
selection (top-1 accuracy).  We also record the mean rank of the LLM's selected tool in
the similarity ranking (chance rank $= 2.5$) and the mean cosine similarity of selected
vs.\ top-similar tools.

\subsection{Positional Bias Analysis}
\label{sec:pos}

For the rotation dataset, we compute the empirical selection rate at each of the four
positions (0--3) and test the null hypothesis of uniform selection (25\% each) using a
chi-square test.  We define \emph{selection stability} as the fraction of base examples
for which the same tool is chosen across all four rotations.  The total variation distance
from the uniform distribution serves as a summary statistic for position preference.

\subsection{Description Perturbation}
\label{sec:perturb}

For each of the three query types, we collect \gptmini selections under all three
conditions (original, generic, adversarial) and record whether the selection changes
relative to the original.  Under the adversarial condition, we further check whether the
model follows the swapped description (selecting the tool that received the correct tool's
description) or the tool name.

\subsection{Layer-wise Probing with \gpttwo}
\label{sec:probing}

We extract hidden states from all 12 transformer layers of \gpttwo (117M parameters,
768-dimensional hidden states) on the 100 probing prompts.  Prompts are formatted to fit
within 512 tokens.  We study two probing targets:

\begin{enumerate}[leftmargin=*,itemsep=0pt,topsep=2pt]
    \item \textbf{Tool category} (5-class): predicts which of the five tool categories is
          relevant for the given query.
    \item \textbf{Correct tool position} (4-class): predicts which of the four tool slots
          contains the semantically correct tool.
\end{enumerate}

For each target we train a logistic regression probe with $L_2$ regularization
($C = 1.0$, max iterations $= 2000$) on each layer's representations, evaluated with
5-fold stratified cross-validation (seed~42).  We study two representation aggregations:
(a)~\emph{mean pooling} over all token positions, and (b)~\emph{last-token}
representation, which corresponds to the position at which the model would generate its
answer in autoregressive decoding.  Chance accuracy is 20\% for the category probe and
25\% for the position probe.

All experiments run on CPU only (no GPU required); LLM experiments take approximately
15 minutes and probing approximately 8 minutes.
